% ============================================
% Assistant Instructional Professor Cover Letter – Harris School of Public Policy
% ============================================

\documentclass[11pt,letterpaper]{letter}

\usepackage{geometry}
\usepackage{microtype}
\usepackage[hidelinks]{hyperref}

\geometry{left=25mm,right=25mm,top=25mm,bottom=25mm}

\signature{Decory Edwards}
\address{
Department of Economics \\
Johns Hopkins University \\
Baltimore, MD 21218 \\
\texttt{dedwar65@jh.edu}
}

\begin{document}

\begin{letter}{Search Committee \\
Harris School of Public Policy \\
University of Chicago \\
Chicago, IL}

\opening{Dear Members of the Search Committee,}

I am writing to apply for the Assistant Instructional Professor position at the Harris School of Public Policy. I am a Ph.D.\ candidate in Economics at Johns Hopkins University, advised by Professors Christopher D.~Carroll, Nicholas W.~Papageorge, and Stelios Fourakis, and I expect to receive my degree in May 2026. My research fields are macroeconomics, wealth and income dynamics, and computational economics.

My research agenda focuses on understanding the determinants of wealth inequality and the mechanisms through which heterogeneity in returns to wealth shapes aggregate outcomes. In my job-market paper, I estimate the degree of heterogeneity in individual portfolio returns using micro-data from the Survey of Consumer Finances and simulate a structural model in which households face idiosyncratic return risk. I show that allowing for persistent heterogeneity in returns substantially improves the model’s ability to match the empirical distribution of wealth, offering a tractable way to connect micro-level return dispersion to macroeconomic inequality.

In a second project, I use the Health and Retirement Study to examine the relationship between interpersonal trust and portfolio performance. Building on the literature that finds a hump-shaped relationship between trust and income, I show that a similar relationship holds for individual returns to wealth measured in the HRS data, highlighting a behavioral channel through which social attitudes shape saving and investment outcomes. This work also provides new estimates of the persistent component of returns to net worth, which motivate the calibration of heterogeneity in my structural model.

The Harris School is an especially strong fit for my interests in teaching applied macroeconomics and finance in a policy-oriented setting. Harris’s curriculum emphasizes analytical rigor, empirical reasoning, and clear connections between economic theory and real-world policy questions. These priorities align closely with how I approach the classroom. I am particularly drawn to the opportunity to teach macroeconomics, international trade, and finance to students who intend to apply economic reasoning directly to public policy, regulation, and institutional design.

\textbf{Teaching.} I have experience teaching and supporting courses in \textit{Principles of Macroeconomics}, \textit{Intermediate Macroeconomic Theory}, and applied macroeconomic policy. My teaching emphasizes economic intuition, careful use of data, and translating formal models into tools that students can use to evaluate policy tradeoffs. At Harris, I would be well prepared to teach courses in macroeconomics, international trade, finance, and financial regulation, with an emphasis on policy relevance and empirical grounding. I value clear communication and structured problem-solving, and I aim to make quantitative material accessible to students with diverse academic backgrounds.

I am enthusiastic about the Harris School’s mission to train students for policy-facing careers and its strong culture of teaching excellence. I would welcome the opportunity to contribute to the instructional mission of the School and to work closely with students in a rigorous, applied learning environment.

Thank you very much for your time and consideration. I would be happy to provide any additional materials or information.

\closing{Sincerely,}

\end{letter}
\end{document}
